\documentclass[../DoAn.tex]{subfiles}
\begin{document}
\chapter{Giải pháp và đóng góp}

Chương này trình bày các giải pháp và đóng góp chính trong quá trình xây dựng hệ thống TheWeekend. Khác với Chương 4 tập trung vào kiến trúc, mô hình dữ liệu, thành phần triển khai và kết quả xây dựng, chương này phân tích các đóng góp ở mức tổ chức nghiệp vụ và chức năng hệ thống. Các nội dung được trình bày theo các nhóm giải pháp chính: tổ chức dữ liệu địa điểm phục vụ tra cứu và gợi ý, xây dựng chatbot dựa trên dữ liệu nội bộ, số hóa quy trình đặt vé và thanh toán, phát hành vé QR và hỗ trợ check-in bằng ứng dụng Android scanner, cùng với phần quản trị dữ liệu và vận hành hệ thống.

Các đóng góp được xem xét trong phạm vi một đồ án xây dựng hệ thống web hoàn chỉnh, không trình bày như kết quả nghiên cứu thuật toán đã được đánh giá định lượng. Hệ thống đã được kiểm thử thủ công trên môi trường local cho các luồng chính, tuy nhiên chưa có khảo sát người dùng, kiểm thử tải hoặc kiểm thử hiệu năng. Vì vậy, chương này không đưa ra kết luận về hiệu năng, mức độ chính xác, độ ổn định khi có nhiều người dùng đồng thời hay mức độ hài lòng của người dùng. Mục tiêu của chương là làm rõ cách TheWeekend giải quyết các vấn đề chức năng đã xác định và những giới hạn còn tồn tại cần tiếp tục hoàn thiện.

\section{Tổ chức hệ thống tra cứu và gợi ý địa điểm dựa trên dữ liệu nội bộ}

Một trong các vấn đề trung tâm của TheWeekend là hỗ trợ người dùng tìm kiếm địa điểm vui chơi cuối tuần phù hợp cho trẻ em và gia đình có trẻ em. Khi lựa chọn một địa điểm, người dùng thường cần quan tâm đồng thời đến nhiều thông tin như vị trí, độ tuổi phù hợp, giá vé, tiện ích, giờ mở cửa, hình ảnh, đánh giá và mức độ phù hợp với nhu cầu của gia đình. Nếu các thông tin này được trình bày rời rạc hoặc nằm ở nhiều nguồn khác nhau, quá trình lựa chọn có thể trở nên mất thời gian và khó so sánh.

Giải pháp của TheWeekend là tổ chức dữ liệu địa điểm trong một hệ thống nội bộ thay vì chỉ phụ thuộc vào việc hiển thị thông tin thô. Dữ liệu địa điểm không chỉ gồm tên và địa chỉ, mà còn bao gồm các thuộc tính mô tả phục vụ nghiệp vụ như danh mục, tag, tiện ích, độ tuổi phù hợp, hình ảnh, đánh giá và thông tin liên quan đến vé nếu có. Việc tập trung dữ liệu này giúp hệ thống có một nguồn thông tin thống nhất để phục vụ nhiều chức năng khác nhau, từ tìm kiếm, xem chi tiết đến gợi ý địa điểm.

Các nhóm dữ liệu địa điểm, danh mục, tag, đánh giá và yêu thích bổ trợ lẫn nhau và cùng phục vụ ba chức năng: tìm kiếm, xem chi tiết và gợi ý. Việc tập trung dữ liệu nội bộ giúp các chức năng này dùng chung một nguồn thông tin thống nhất, tạo ra luồng sử dụng liên tục thay vì tách rời từng thao tác. Riêng cơ chế lọc và xếp hạng gợi ý theo tiêu chí (vị trí, độ tuổi, thẻ, danh mục, giá vé, đánh giá) đã được trình bày chi tiết ở Mục~\ref{sec:recommendation}.

Đóng góp của nhóm chức năng này nằm ở việc hình thành một luồng tra cứu tương đối hoàn chỉnh: người dùng có thể bắt đầu từ tìm kiếm địa điểm, mở trang chi tiết để xem thông tin, tham khảo đánh giá, lưu địa điểm yêu thích và tiếp tục nhận gợi ý từ dữ liệu hệ thống. Luồng này phù hợp với nhu cầu tham khảo và lựa chọn địa điểm vui chơi, đặc biệt trong bối cảnh người dùng cần cân nhắc nhiều tiêu chí trước khi quyết định.

Việc tổ chức dữ liệu nội bộ cũng hỗ trợ phía quản trị trong quá trình cập nhật và kiểm soát dữ liệu địa điểm. Khi dữ liệu địa điểm, danh mục và tag được quản lý trong hệ thống, quản trị viên có thể bổ sung hoặc điều chỉnh thông tin theo nhu cầu vận hành. Điều này giúp dữ liệu phục vụ người dùng không hoàn toàn phụ thuộc vào một nguồn bên ngoài, đồng thời tạo điều kiện để hệ thống mở rộng thêm các tiêu chí tra cứu trong tương lai.

Giới hạn của nhóm giải pháp này là chưa có dữ liệu đánh giá chính thức về chất lượng gợi ý địa điểm. Báo cáo không khẳng định cơ chế gợi ý đạt độ chính xác cao hay tối ưu, vì chưa có bộ dữ liệu kiểm thử, tiêu chí đánh giá định lượng hoặc phản hồi người dùng chính thức. Vì vậy, trong phạm vi đồ án, chức năng này được mô tả là cơ chế gợi ý dựa trên dữ liệu nội bộ, không trình bày như một thuật toán gợi ý đã được kiểm chứng về hiệu quả.

\section{Chatbot hỏi đáp dựa trên Gemini API và dữ liệu MongoDB}

Bên cạnh nhu cầu tra cứu bằng danh sách và bộ lọc, người dùng có thể muốn đặt câu hỏi trực tiếp bằng ngôn ngữ tự nhiên. Ví dụ, người dùng có thể muốn hỏi về địa điểm phù hợp cho trẻ nhỏ, địa điểm có hoạt động cụ thể, hoặc thông tin liên quan đến tiện ích và giá vé. Trong các tình huống này, giao diện tìm kiếm truyền thống có thể chưa đủ thuận tiện nếu người dùng chưa biết rõ nên chọn tiêu chí nào hoặc phải đọc nhiều trang chi tiết để tìm câu trả lời.

TheWeekend bổ sung chatbot nhằm hỗ trợ người dùng hỏi đáp nhanh trong phạm vi dữ liệu của hệ thống. Chức năng này không thay thế hoàn toàn chức năng tìm kiếm và xem chi tiết địa điểm, mà đóng vai trò như một kênh hỗ trợ bổ sung. Người dùng có thể đặt câu hỏi tự nhiên, sau đó hệ thống xử lý câu hỏi, truy hồi dữ liệu liên quan từ MongoDB và sử dụng Gemini API để sinh phản hồi dựa trên ngữ cảnh đã chuẩn bị.

Cách tiếp cận này kết hợp dữ liệu có cấu trúc trong MongoDB với khả năng diễn đạt của mô hình ngôn ngữ, đồng thời giới hạn phạm vi trả lời trong dữ liệu nội bộ để chatbot không trả lời vượt ra ngoài phạm vi ứng dụng. Nhờ đó, người dùng tiếp cận thông tin địa điểm theo cách gần với hội thoại hơn so với việc chỉ đọc danh sách kết quả. Luồng xử lý chi tiết (truy hồi dữ liệu, tạo ngữ cảnh và lời nhắc, giới hạn phạm vi, xử lý khi không có dữ liệu phù hợp) đã được trình bày ở Mục~\ref{sec:chatbot-mechanism}.

Đóng góp của chức năng chatbot nằm ở việc bổ sung một lớp tương tác mới giữa người dùng và dữ liệu địa điểm. Thay vì chỉ dựa vào thao tác nhập từ khóa, chọn bộ lọc hoặc mở từng trang chi tiết, người dùng có thể hỏi trực tiếp về nhu cầu của mình. Điều này có ý nghĩa trong các tình huống người dùng chưa xác định rõ tiêu chí hoặc muốn được hỗ trợ tổng hợp thông tin nhanh từ dữ liệu có sẵn.

Chatbot cũng giúp kết nối các nhóm dữ liệu khác nhau trong hệ thống. Dữ liệu địa điểm cung cấp thông tin mô tả, dữ liệu tag hỗ trợ xác định đặc điểm, còn dữ liệu đánh giá có thể bổ sung thông tin tham khảo từ người dùng. Khi được sử dụng làm ngữ cảnh cho chatbot, các nhóm dữ liệu này không chỉ phục vụ hiển thị trên giao diện chi tiết mà còn hỗ trợ quá trình hỏi đáp bằng ngôn ngữ tự nhiên.

Giới hạn của chức năng này là chưa có đánh giá định lượng về chất lượng câu trả lời. Báo cáo cũng chưa có bộ câu hỏi chuẩn để kiểm thử mức độ phù hợp, tính nhất quán hoặc độ bao phủ của phản hồi chatbot. Vì vậy, đóng góp của chatbot trong đồ án được xác định ở mức tích hợp chức năng hỏi đáp dựa trên Gemini API và dữ liệu MongoDB, chưa khẳng định độ chính xác hay hiệu quả sử dụng trong điều kiện thực tế.

\section{Số hóa quy trình đặt vé, thanh toán và phát hành vé QR}

Đối với các địa điểm có bán vé, trải nghiệm của người dùng không chỉ dừng lại ở việc tra cứu thông tin. Người dùng còn cần chọn loại vé, tạo đơn đặt vé, thanh toán, nhận vé và xuất trình vé tại điểm đến. Nếu các bước này được xử lý rời rạc, người dùng có thể gặp khó khăn trong việc theo dõi trạng thái đơn, còn phía vận hành khó kiểm tra tính hợp lệ của vé và trạng thái thanh toán.

TheWeekend giải quyết bài toán này bằng cách số hóa quy trình đặt vé và thanh toán trong cùng một hệ thống. Các thành phần nghiệp vụ như booking, payment, payment log, ticket và ticket type được tách thành các nhóm dữ liệu riêng, giúp hệ thống biểu diễn rõ hơn từng phần của quy trình. Booking thể hiện đơn đặt vé, payment thể hiện giao dịch thanh toán, payment log ghi nhận sự kiện thanh toán, ticket type mô tả loại vé, còn ticket đại diện cho vé điện tử được phát hành.

Cách tách các thành phần này giúp luồng đặt vé dễ theo dõi hơn ở mức dữ liệu và nghiệp vụ, đồng thời tránh trộn lẫn dữ liệu đặt vé với dữ liệu thanh toán trong một đối tượng duy nhất. Chuỗi xử lý từ tạo booking, thanh toán PayOS đến phát hành vé QR sau khi giao dịch hợp lệ đã được mô tả ở Chương 4.

PayOS đóng vai trò là dịch vụ thanh toán được tích hợp trong quy trình. Trong phạm vi đồ án, PayOS được sử dụng để tạo yêu cầu thanh toán và hỗ trợ ghi nhận trạng thái giao dịch. Hệ thống không tự xây dựng cổng thanh toán riêng, mà tích hợp dịch vụ ngoài để xử lý phần thanh toán trực tuyến. Điều này phù hợp với mục tiêu của đồ án là xây dựng hệ thống ứng dụng và tích hợp nghiệp vụ, thay vì triển khai hạ tầng thanh toán từ đầu.

Vé QR là thành phần kết nối giữa quá trình đặt vé trực tuyến và quá trình kiểm tra vé tại điểm đến. Sau khi thanh toán hợp lệ, vé điện tử được phát hành với thông tin định danh để người dùng có thể xuất trình. Nhân viên hoặc quản trị viên có thể sử dụng ứng dụng scanner để quét mã QR, gửi dữ liệu về backend và kiểm tra trạng thái vé. Như vậy, vé QR không chỉ là hình thức hiển thị cho người dùng mà còn là cầu nối giữa frontend, backend và ứng dụng Android scanner.

Đóng góp của nhóm chức năng này là thống nhất luồng từ chọn vé, tạo booking, thanh toán, phát hành vé QR đến theo dõi vé trong cùng một hệ thống. Đối với người dùng, luồng này giúp theo dõi quá trình đặt vé và nhận vé rõ ràng hơn. Đối với quản trị viên, dữ liệu booking, payment và ticket hỗ trợ việc quan sát nghiệp vụ vé và thanh toán. Đối với nhân viên, vé QR cung cấp đầu vào cho quy trình xác minh và check-in tại điểm đến.

Việc số hóa quy trình vé cũng giúp giảm sự phụ thuộc vào kiểm tra thủ công. Khi vé được phát hành dưới dạng điện tử và có mã QR, hệ thống có thể xác minh vé dựa trên dữ liệu backend thay vì chỉ dựa trên thông tin người dùng cung cấp bằng lời nói hoặc mã văn bản rời rạc. Điều này tạo tiền đề cho quy trình vận hành nhất quán hơn giữa lúc người dùng đặt vé và lúc nhân viên kiểm tra vé.

Tuy nhiên, báo cáo chưa đủ dữ liệu để kết luận chính sách hủy vé, hoàn tiền, hết hạn thanh toán hoặc vòng đời đầy đủ của vé như một quy trình nghiệp vụ chính thức. Một số trạng thái có thể xuất hiện trong mô hình dữ liệu hoặc mã nguồn, nhưng việc mô tả thành chính sách vận hành cần có yêu cầu nghiệp vụ rõ ràng hơn. Vì vậy, chương này chỉ phân tích đóng góp của việc số hóa luồng đặt vé, thanh toán và phát hành vé QR, không tự bổ sung quy trình hủy vé hoặc hoàn tiền.

\section{Ứng dụng Android scanner cho nghiệp vụ check-in}

Sau khi vé QR được phát hành, hệ thống cần một cơ chế để kiểm tra vé tại điểm đến. Nếu nhân viên chỉ kiểm tra thủ công bằng cách đọc mã vé hoặc đối chiếu thông tin trên màn hình người dùng, quá trình này có thể thiếu nhất quán và khó cập nhật trạng thái sử dụng vé kịp thời. Vì vậy, TheWeekend bổ sung ứng dụng Android scanner để hỗ trợ nghiệp vụ xác minh và check-in vé.

Việc sử dụng ứng dụng Android scanner phù hợp với bối cảnh nhân viên thao tác tại quầy hoặc tại cổng vào. Thiết bị Android có thể sử dụng camera để quét mã QR và gửi dữ liệu về backend. So với việc kiểm tra thủ công, cách tiếp cận này giúp thao tác xác minh vé có quy trình rõ ràng hơn: quét mã, nhận thông tin xác minh, sau đó xác nhận check-in nếu vé hợp lệ.

Ứng dụng scanner được xây dựng bằng Android native Java, dùng ZXing để đọc mã QR và OkHttp để gọi backend; toàn bộ logic kiểm tra vé được xử lý tập trung ở backend, còn ứng dụng chỉ tập trung vào thao tác tại thiết bị. Cách scanner xác thực với backend, gọi API kiểm tra vé, hiển thị kết quả xác minh và xử lý khi mất mạng đã được trình bày ở Mục~\ref{sec:checkin-mechanism}.

Đóng góp của ứng dụng scanner là kết nối vé điện tử với vận hành tại điểm đến. Vé QR được phát hành sau thanh toán sẽ không chỉ nằm trong giao diện người dùng, mà còn được sử dụng trong quy trình kiểm tra thực tế. Nhờ đó, hệ thống tạo ra một luồng liên tục từ đặt vé trực tuyến đến xác minh vé tại điểm đến, thay vì dừng lại ở việc hiển thị vé cho người dùng.

Đối với nhân viên hoặc quản trị viên, scanner giúp tách rõ thao tác vận hành vé khỏi giao diện người dùng thông thường. Nhân viên không cần truy cập toàn bộ hệ thống quản trị để thực hiện việc quét QR, mà có một ứng dụng tập trung vào nhiệm vụ check-in. Điều này phù hợp với nguyên tắc phân quyền theo vai trò, trong đó nhân viên chỉ cần sử dụng các chức năng phục vụ nghiệp vụ tại điểm đến.

Trong phạm vi đồ án, ứng dụng Android scanner đã được kiểm thử thủ công trên thiết bị Android thật trong môi trường local, bao gồm đăng nhập, quét QR, xác minh vé và ghi nhận check-in. Tuy nhiên, báo cáo chưa có kiểm thử thực địa tại điểm đến, chưa đánh giá tốc độ quét, độ ổn định mạng, khả năng hoạt động trên nhiều thiết bị Android hoặc độ chính xác trong các điều kiện ánh sáng khác nhau. Do đó, đóng góp của scanner được trình bày ở mức xây dựng chức năng và xác nhận luồng nghiệp vụ trong môi trường local, chưa kết luận về hiệu quả vận hành thực tế.

\section{Tổ chức quản trị dữ liệu và vận hành hệ thống}

Một hệ thống gợi ý và đặt vé địa điểm không chỉ cần giao diện cho người dùng cuối, mà còn cần công cụ quản trị dữ liệu tương ứng. Nếu dữ liệu địa điểm, đánh giá, ảnh đánh giá, vé và thanh toán không được quản lý tập trung, hệ thống khó duy trì tính nhất quán trong quá trình sử dụng. Vì vậy, phần quản trị là một đóng góp quan trọng của TheWeekend bên cạnh các chức năng dành cho người dùng.

Quản trị viên trong hệ thống có vai trò quản lý các dữ liệu vận hành chính, bao gồm người dùng, địa điểm, tag, đánh giá, ảnh đánh giá, vé, thanh toán và thông báo. Các nhóm dữ liệu này liên quan trực tiếp đến trải nghiệm của người dùng cuối. Ví dụ, dữ liệu địa điểm ảnh hưởng đến kết quả tìm kiếm và gợi ý; dữ liệu đánh giá ảnh hưởng đến thông tin tham khảo; dữ liệu vé và thanh toán ảnh hưởng đến nghiệp vụ đặt vé.

Việc phân quyền admin và staff giúp tách biệt giữa quản trị hệ thống và vận hành check-in. Quản trị viên có phạm vi thao tác rộng hơn, bao gồm quản lý dữ liệu và theo dõi nghiệp vụ. Nhân viên tập trung vào các thao tác liên quan đến xác minh và check-in vé. Cách phân quyền này giúp hệ thống phản ánh tốt hơn các vai trò sử dụng thực tế, đồng thời hạn chế việc một tài khoản vận hành tại điểm đến phải có toàn bộ quyền quản trị.

Chức năng duyệt ảnh đánh giá có ý nghĩa trong việc kiểm soát nội dung do người dùng gửi lên. Ảnh đánh giá có thể giúp thông tin địa điểm trực quan hơn, nhưng cũng cần được xử lý trước khi hiển thị công khai để tránh nội dung không phù hợp. Việc bổ sung trạng thái xử lý ảnh đánh giá giúp quản trị viên có cơ sở để duyệt hoặc từ chối ảnh, thay vì đưa trực tiếp mọi nội dung người dùng gửi vào hệ thống.

Quản lý vé và thanh toán hỗ trợ phía vận hành theo dõi các booking, giao dịch và vé đã phát hành. Nhóm chức năng này không chỉ phục vụ người dùng trong quá trình mua vé, mà còn giúp quản trị viên quan sát trạng thái nghiệp vụ sau khi đơn được tạo. Khi kết hợp với ứng dụng scanner, dữ liệu vé có thể tiếp tục được sử dụng trong bước xác minh và check-in tại điểm đến.

Thông báo trong hệ thống giúp ghi nhận và truyền đạt một số trạng thái hoặc sự kiện liên quan đến người dùng. Trong phạm vi đồ án, thông báo được xem như một thành phần hỗ trợ trải nghiệm và vận hành, giúp người dùng nhận biết các cập nhật liên quan tới tài khoản hoặc nghiệp vụ. Tuy nhiên, báo cáo chưa trình bày thông báo như một hệ thống giám sát hay kênh thông báo thời gian thực đã được đánh giá.

Đóng góp của nhóm quản trị là giúp TheWeekend không chỉ là một giao diện tra cứu địa điểm, mà còn có phần quản lý và vận hành dữ liệu tương ứng. Nhờ đó, các chức năng người dùng, vé, thanh toán, đánh giá và check-in có thể được đặt trong một hệ thống có phân quyền và có dữ liệu quản trị đi kèm. Điều này làm cho phạm vi đồ án hoàn chỉnh hơn so với một ứng dụng chỉ hiển thị danh sách địa điểm.

Giới hạn của nhóm chức năng này là chưa có dashboard thống kê chính thức và chưa có dữ liệu đánh giá hiệu quả quản trị. Báo cáo cũng chưa có số liệu về số lượng bản ghi, số lượng giao dịch hoặc tần suất sử dụng với dữ liệu thật. Vì vậy, các đóng góp quản trị được mô tả ở mức chức năng đã xây dựng và tổ chức nghiệp vụ, chưa khẳng định hiệu quả khi sử dụng ở quy mô lớn.

\section{Tổng kết chương}

Chương này đã trình bày các nhóm giải pháp và đóng góp chính của TheWeekend. Các đóng góp bao gồm tổ chức dữ liệu địa điểm phục vụ tra cứu và gợi ý, xây dựng chatbot dựa trên Gemini API và dữ liệu MongoDB, số hóa quy trình đặt vé và thanh toán, phát hành vé QR và check-in bằng Android scanner, cùng với tổ chức phần quản trị dữ liệu và vận hành hệ thống. Các nhóm đóng góp này liên kết với nhau để tạo thành một hệ thống hỗ trợ người dùng từ bước tìm kiếm địa điểm đến bước đặt vé và sử dụng vé tại điểm đến.

Các đóng góp trong chương được xác định ở mức xây dựng hệ thống và tổ chức nghiệp vụ, không phải kết quả nghiên cứu thuật toán định lượng. Những hạn chế còn tồn tại — thiếu khảo sát người dùng, kiểm thử định lượng và dữ liệu nghiệp vụ cho một số chính sách — được tổng hợp cùng các hướng phát triển trong Chương 6, nhằm tránh phóng đại phạm vi của đồ án.

\end{document}
